La evolución del spread bancario en moneda nacional (MN) y moneda extranjera (ME) presenta diferencias importantes durante el período analizado. Tomando como referencia el primer trimestre de 2015 (índice base 100), el spread bancario en moneda nacional muestra una tendencia descendente y relativamente estable, con una variación aproximada entre 100.5 y 92.3. Por su parte, el spread bancario en moneda extranjera evidencia una mayor amplitud en sus fluctuaciones, alcanzando valores entre 113.9 y 97.4, lo que refleja una mayor variabilidad respecto a su período base.
A partir de 2019 el spread bancario en moneda extranjera incrementa su volatilidad, comportamiento que comienza a moderarse hacia 2023. En contraste, el spread en moneda nacional mantiene un comportamiento más estable, aunque desde finales de 2021 inicia una reducción gradual que lo lleva a registrar su nivel mínimo durante el último trimestre de 2024. Estas variaciones coinciden temporalmente con un período de elevada incertidumbre económica internacional, caracterizado por los efectos de la pandemia de COVID-19 y posteriormente por el conflicto entre Rusia y Ucrania, eventos que influyeron sobre las condiciones financieras y los mercados internacionales.
En conjunto, los resultados sugieren que el spread bancario en moneda nacional ha mantenido un comportamiento relativamente más estable que el observado en moneda extranjera. Esta diferencia puede interpretarse como un reflejo de la mayor sensibilidad de las operaciones denominadas en moneda extranjera frente a factores externos, mientras que el mercado en moneda nacional tiende a responder principalmente a las condiciones monetarias y financieras internas. Por ello, el seguimiento conjunto de ambos indicadores constituye una herramienta relevante para evaluar la estabilidad del sistema financiero y la evolución de las condiciones de financiamiento en el país.